
\begin{titlepage}
    \centering

    {\largefont \textbf{ABSTRACT}}
    
    \begin{justify}
      \normalsizefont{
With the development of mobile computing, projects are increasingly designed with a focus on the qualities of this type of system. In contrast to web applications, the offline state is not considered an error or a temporary condition, but rather a fact that mobile development must take into account. Thus, the Offline-First paradigm has been consolidated as a reality. This paradigm describes, among other aspects, the need for data synchronization across devices connected to the system, in order to maintain the consistency and availability of such data. In scenarios where processes are centralized, such as on a web server, there are specific strategies for maintaining files and states; the same applies to decentralized computing. Peer-to-Peer networks are a good alternative to servers due to their low cost and privacy, and they have gained further traction with the rise of blockchains. However, Peer-to-Peer networks present their own synchronization challenges: peer churn, replica consistency, and synchronization efficiency. Therefore, studying synchronization strategies in Peer-to-Peer networks within Offline-First systems is essential to guide decision-making in selecting data synchronization implementations that are congruent with the system in production. This study must be carried out using metrics in a controlled environment with nodes operating in a decentralized network, in order to collect data suitable for analysis and comparison of the advantages, limitations, and behaviors of different synchronization strategies. In the same sense, the development of a real Offline-First application to provide substance to each of these strategies is indispensable for validating the discussions conducted about each synchronization approach.
} \\ [1cm]
      \normalsizefont \textbf{Keywords: } {Offline-First, Peer-to-Peer, Data Synchronization, CRDT, CAP.}
    \end{justify}

    \vfill
\end{titlepage}
